% !TeX root = ../main.tex
\section{SaiLoR from the Reviewer's View}\label{sec:tool}

In this section, we walk through \SaiLoR along the activities of an \SLR.
We use the running example from \autoref{sec:approach-example}.
% TODO: All figures show the user interface before the move to Tauri (user guide, commit 1c0ec0f). Retake them on the 2.0.0 build, see figures.tex.

\begin{figure*}[t]
    \centering
    \includegraphics[width=0.85\textwidth]{\shotAnnotate}
    \caption{Main view of \SaiLoR with the paper list on the left, the PDF in the middle, and the extraction form of the running example for reviewer seat 1 on the right}
    \label{fig:annotate}
\end{figure*}

\subsection{Planning: Protocol and Schema}\label{sec:tool-planning}

% Project editor
Reviewers start in the project editor.
There, they set the project title, the number of reviewer seats, and whether the project is a screening project.
An optional review protocol records the research questions, search strings, databases, search date, and further notes.
The main part of the editor is the schema, as shown in \autoref{fig:schema}.
Each row defines one node with its name, type, minimum and maximum count, options, description, and visibility condition.
In the figure, the field \emph{Exclusion Reason} is only visible if the field \emph{Relevant} has the answer \enquote{No}.
Reviewers can reorder and nest rows by drag and drop.

\begin{figure}[t]
    \centering
    \includegraphics[width=\columnwidth]{\shotSchema}
    \caption{Schema editor with field types, counts, options, and visibility conditions}
    \label{fig:schema}
\end{figure}

\subsection{Conducting: Import and Screening}\label{sec:tool-screening}

% Import and duplicates
Reviewers import references from BibTeX, RIS, or CSL-JSON files, and they can import a folder of PDFs including its subfolders.
During import, \SaiLoR sorts each record into one of three classes: new, certain duplicate, or probable duplicate.
A certain duplicate has the same normalized DOI or title as a paper in the project.
A probable duplicate has a fuzzy title similarity of at least $0.90$, or the same base title and an author similarity of at least $0.50$.
The reviewer must decide on each probable duplicate.
If the reviewer merges two records, \SaiLoR only fills fields that are empty in the existing paper.

% Screening
In a screening project, each reviewer decides per paper whether to include it, to exclude it with a reason, or to leave it undecided.
\autoref{fig:screening} shows the screening view of the running example.
The view shows the abstract instead of the PDF, because reviewers usually screen titles and abstracts first.
If the abstract was extracted from the PDF and not imported, \SaiLoR shows a warning, because the extraction can be wrong.
Reviewers can screen with the keyboard: \texttt{I} includes, \texttt{E} excludes, \texttt{U} clears the decision, and the number keys exclude with the matching reason.
After a decision, \SaiLoR moves on to the next paper.
A progress line shows how many papers the reviewer screened, included, and excluded.
The screening summary lists these counts with one row per exclusion reason, which is the input for a \PRISMA flow diagram~\cite{page2021prisma}.
Validation reports excluded papers without a reason.

\begin{figure}[t]
    \centering
    \includegraphics[width=\columnwidth]{\shotScreening}
    \caption{Screening view with the abstract, the decision buttons, and the progress line}
    \label{fig:screening}
\end{figure}

% Carry over
When screening is done, reviewers create the extraction project from the screening project.
\SaiLoR carries over all included papers, never carries over excluded papers, and lets the reviewer choose whether to carry over undecided papers.
The new project records where it came from and how many papers were included, excluded, and carried over.

\subsection{Conducting: Data Extraction}\label{sec:tool-extraction}

% Extraction view
\autoref{fig:annotate} shows the main view during extraction.
The left pane lists the papers, the middle pane shows the PDF, and the right pane shows the form that \SaiLoR generates from the schema.
Repeatable groups such as \emph{Findings} get buttons to add and remove entries.
A dot next to each paper shows how complete the paper is for the current seat.
Reviewers mark a paper as finished explicitly.
Thus, \enquote{finished} is a decision of the reviewer and not just the absence of empty fields.

% Evidence
Reviewers can copy selected text from the PDF into a field.
They can also highlight passages in five colors and add notes.
Each reviewer seat has its own marks.
A reviewer can link a mark to one or more fields as evidence.
The link popover of a field lists the linked marks, and a click on a mark jumps to the passage in the PDF.
In this way, the project stores where in the paper an answer comes from.
Reviewers can export their marks as standard PDF annotations into a new PDF file.

\subsection{Conducting: Consolidation}\label{sec:tool-consolidation}

% Seats
In a project with several reviewers, every reviewer must pick a seat before editing.
A reviewer only sees the answers of their own seat.
The consolidator works in the additional consolidation seat.

% Alignment
Before \SaiLoR can compare answers, it must find out which entries of a repeatable group belong together.
For example, two reviewers can list the findings of a paper in different orders.
\SaiLoR matches the entries with a maximum-weight assignment over their pairwise similarity, and it keeps entries with a similarity of $0.5$ or lower apart.
The alignment is deterministic, so that all reviewers see the same result.
Before \SaiLoR compares two answers, it trims them, converts them to lowercase, and collapses whitespace.
If at least two reviewers answered a field and all answers are equal, \SaiLoR fills the consolidated value automatically and marks it as auto-filled.

% Compare and adopt
For the remaining fields, the consolidator opens a comparison, as shown in \autoref{fig:compare}.
The comparison lists the answers of all reviewers.
The consolidator adopts one answer or marks the answers as equivalent.
The disagreements of a paper or of the whole project can be exported as text, for example, to discuss them in a meeting.

\begin{figure}[t]
    \centering
    \includegraphics[width=\columnwidth]{\shotCompare}
    \caption{Comparison of two reviewers' answers in the consolidation seat}
    \label{fig:compare}
\end{figure}

% Agreement
\SaiLoR also computes the inter-rater agreement.
For exactly two reviewers, it computes Cohen's $\kappa$~\cite{cohen1960coefficient}, and for two or more reviewers Fleiss' $\kappa$~\cite{fleiss1971measuring}.
In both cases, it also computes Krippendorff's $\alpha$ for nominal data~\cite{krippendorff2018content}, which tolerates missing answers.
\SaiLoR shows the agreement overall and per field.
If a statistic is not defined, for example, because all reviewers gave the same answer to every item, \SaiLoR reports the case instead of a number.
In a screening project, the agreement only considers the decision.

\subsection{Collaboration with Git}\label{sec:tool-git}

% Commit and pull
\SaiLoR has a built-in git view that uses the git installation of the user.
Instead of a line diff of JSON files, the view lists each changed field with its previous and current value.
Reviewers select which field changes they commit.
A pull fetches the upstream branch and reports whether the project is up to date, can be fast-forwarded, or needs a merge.
\SaiLoR only pushes when the reviewer asks for it.
It also supports branches, a project-scoped history, and stashes of the project's own files.

% Merge
For a merge, \SaiLoR does not use the line-based merge of git.
It runs a three-way merge on the level of fields.
If only one side changed a field, the merge keeps that change.
The schema is merged node by node, and answers under a removed field are kept.
Only if both sides changed the same field to different values does the reviewer have to decide.
\autoref{fig:merge} shows the dialog for these conflicts.
If a merge leaves conflicts in files that do not belong to the project, \SaiLoR aborts it and restores the previous state.

\begin{figure}[t]
    \centering
    \includegraphics[width=\columnwidth]{\shotMerge}
    \caption{Field-level resolution of merge conflicts after a pull}
    \label{fig:merge}
\end{figure}

\subsection{\acs{LLM} Assistance}\label{sec:tool-llm}

% Providers and opt-in
\LLM assistance is off unless the project enables it.
Reviewers can connect nine provider \APIs, among them Anthropic, OpenAI, Google Gemini, and any OpenAI-compatible server.
Local models run, for example, in Ollama or a llama.cpp server~\cite{llamacpp}.
\SaiLoR stores \API keys in the keychain of the operating system.
Every \LLM window shows a notice that the feature is a preview and that each proposal must be checked.

% Extraction proposals
For extraction, the model reads the paper and proposes answers.
\SaiLoR validates every proposal against the schema.
A review table shows each proposal with the field, the value, a quote from the paper, and a confidence.
In agent mode, the model can search and read pages of the paper, query OpenAlex and Crossref, and fetch web pages.
\SaiLoR checks whether the quoted evidence occurs in the paper, and an optional judge model checks the proposals in a bounded loop.
Nothing changes in the project until the reviewer applies the selected rows.
Applying is one undo step and does not overwrite existing answers.
\SaiLoR marks each field that the model wrote until the reviewer interacts with the field.
For each apply, \SaiLoR records the provider, the model, and the time in the project, and the paper shows a disclosure line such as \enquote{AI-assisted (2$\times$)}.

% AI seat and screening
In a project with several reviewers, the last seat can be dedicated to the model.
The answers of this seat go through the same consolidation as the answers of the human reviewers, so that the agreement statistics also cover the model.
For screening, the model proposes a decision and a reason per paper.
\SaiLoR never pre-selects an exclusion, and it pre-selects an inclusion only if the confidence is at least $0.8$.
